\documentclass[10pt,a4paper]{amsart}
\usepackage[margin=19mm]{geometry}
\usepackage{amsmath,amssymb,mathtools}
\usepackage[scr=rsfs,cal=euler]{mathalfa}
\usepackage{xfrac}
\usepackage[unicode,hidelinks,bookmarks=false]{hyperref}
\newcommand{\ie}{i.e.\ }
\newcommand{\cf}{cf.\ }
\newcommand{\resp}{respectively\ }
\newcommand{\Subsection}{\subsection}
\providecommand{\nobreakdash}{\nobreak\hbox{-}}
\setlength{\emergencystretch}{2em}
\multlinegap0pt
\usepackage{bm}
\usepackage{bbm}
\usepackage{dsfont}
\usepackage{xspace}
\usepackage{cancel}
\usepackage{mathtools}
\usepackage{xcolor}
\usepackage{amsthm}
\makeatletter
\@ifundefined{definition}{\newtheorem{definition}{Definition}}{}
\@ifundefined{corollary}{\newtheorem{corollary}[definition]{Corollary}}{}
\@ifundefined{lemma}{\newtheorem{lemma}[definition]{Lemma}}{}
\@ifundefined{theorem}{\newtheorem{theorem}[definition]{Theorem}}{}
\makeatother
\def \N {\mathbb{N}}
\def \R {\mathbb{R}}
\def \XX {\mathcal{X}}
\title[Regularizing effects of absorption terms in local-nonlocal m]{Regularizing effects of absorption terms in local-nonlocal mild singular problems}
\author{Stefano Biagi, Enzo Maria Merlino, Eugenio Vecchi}
\date{Source: arXiv:2506.11656v1 (CC BY 4.0)}
\begin{document}
\maketitle

\noindent\emph{Source and licence.} Regularizing effects of absorption terms in local-nonlocal mild singular problems. Version arXiv:2506.11656v1 (CC BY 4.0), distributed under the Creative Commons Attribution 4.0 International licence, as recorded in the arXiv licence metadata for this exact version and shown on the abstract page https://arxiv.org/abs/2506.11656v1. Full rights notice: licenses/arxiv-2506.11656-CC-BY-4.0.txt.\par\smallskip

\noindent\textit{Editorial scope.} Counted unit: the Theorem whose source label is \texttt{thm:pointwise\_talenti} in the article (source0.tex lines 2051--2069, source label \texttt{thm:pointwise\_talenti}), delivered together with the article's own complete original proof of it (source0.tex lines 2070--2240, one \texttt{proof} environment of 171 lines). No mathematics is written, repaired or re-derived: statement and proof are the article's own text line for line. Required context retained uncounted: eq:ModelProblem (lines 226--232); cor:ModelProblemEU (lines 280--289); eq:def\_Tk (lines 376--379); eq:Approx\_Problem (lines 1091--1096); lem:First\_Approx (lines 1105--1112). Every definition, display and label of those lines is retained unchanged. Omitted from this package and not redistributed: 1--225, 233--279, 290--375, 380--1090, 1097--1104, 1113--2050, 2241--3106. Nothing inside a retained excerpt is omitted or altered: every retained body is delivered exactly as the article's lines are written, so no omission receipt of any kind accompanies this package. Citations: the retained excerpts carry no citation command at all, so no bibliography entry of the article is transcribed and no reference is redistributed. Reference adaptation: none was needed and none was made; every reference inside the delivered unit resolves inside it, so no unresolved reference or citation is produced. Printed slips kept literal: no printed word, symbol, hypothesis or number of the counted statement or of its proof was altered. Layout and class: the article's own class is \texttt{amsart} with packages including amsmath, amsfonts, graphicx, framed, amssymb, cancel, xcolor, esint, comment, etex, latexsym, babel, inputenc, amsthm; the wrapper is the lane's pinned \texttt{amsart} editorial wrapper at 10pt on A4, which restates the theorem-like environments the retained text needs and adds the article's own macro definitions that the retained text needs (\texttt{\textbackslash N}, \texttt{\textbackslash R}, \texttt{\textbackslash XX}), with extra wrapper packages (bm, bbm, dsfont, xspace, cancel, mathtools, xcolor). The delivered numbering is the unit's own. Assets: no retained span contains a figure environment, a graphics-inclusion command, a PDF-page inclusion, a table or a TikZ picture; no image file is copied into this package, the package declares no asset dependency and no image file is redistributed with it. Licence: version arXiv:2506.11656v1 of this article is distributed under the Creative Commons Attribution 4.0 International licence, as recorded in the arXiv licence metadata for this exact version and shown on the abstract page https://arxiv.org/abs/2506.11656v1. A full rights notice is delivered at licenses/arxiv-2506.11656-CC-BY-4.0.txt. Characters: every retained excerpt is pure ASCII, so the delivered engine, XeLaTeX with its default Latin Modern text font, reports no missing character.
\begin{equation}\label{eq:ModelProblem}
\left\{ \begin{array}{rl}
-\Delta u + (-\Delta)^s u + u^q = \dfrac{f(x)}{u^{\gamma}} & \textrm{ in } \Omega,\\
u>0 & \textrm{ in } \Omega,\\
u=0 & \textrm{ in } \mathbb{R}^{n}\setminus \Omega,
\end{array}\right.
\end{equation}

\begin{corollary} \label{cor:ModelProblemEU}
	Let $0\leq \gamma\leq 1$, and let $f$ satisfy
	\emph{(H)}$_f$ (with $\theta = \gamma$). According to assumption \emph{(H)}$_g$,  let $q\in\mathbb{R}$ be such that
	$$q\geq \max\{0,q_{\gamma,m}\}, \quad\text{where $q_{\gamma,m}
	= \frac{1-m\gamma}{m-1}$}$$
	(with the convention that $q_{m,\gamma} = -1$ if $\gamma = 1$).
	
	Then, there exists a unique solution $u$ of problem \eqref{eq:ModelProblem} such that
	$$u|_\Omega\in H_0^1(\Omega)\quad\text{and}\quad u\in L^{q+1}(\Omega).$$
\end{corollary}

\begin{align}
& T_{k}(s):= \max \{ -k, \min\{s,k\}\} \quad (s \in \mathbb{R}) \label{eq:def_Tk} \\
& G_{k}(s):= (|s|-k)^{+} \, \mathrm{sign}(s) \quad (s\in \mathbb{R}) \label{eq:def_Gk}
\end{align}

\begin{equation}\label{eq:Approx_Problem}
\left\{ \begin{array}{rl}
\mathcal{L}u_{n,k} + g_{k}(u_{n,k}) = h_n(u_{n,k})f_n & \textrm{in } \Omega,\\
u_{n,k} \geq 0 & \textrm{in } \Omega,\\
u_{n,k} = 0 & \textrm{in } \mathbb{R}^{n}\setminus \Omega.\end{array}\right.
\end{equation}  

\begin{lemma}\label{lem:First_Approx}
The approximating problem \eqref{eq:Approx_Problem} admits a weak solution $u_{n,k}\in \mathcal{X}^{1,2}(\Omega)$ such that
\begin{itemize}
\item[i)] $u_{n,k}\in L^{\infty}(\Omega)$;
\item[ii)] $u_{n,k}$ is bounded in $L^{\infty}(\Omega)$ uniformly in $k\in\mathbb{N}$;
\item[iii)] $u_{n,k}$ is bounded in $\mathcal{X}^{1,2}(\Omega)$ uniformly in $k\in\mathbb{N}$.
\end{itemize}
\end{lemma}  

\begin{theorem}\label{thm:pointwise_talenti}
Let the assumptions of Corollary \ref{cor:ModelProblemEU}, and let $u$ be the unique solu\-tion to the model problem 
\begin{equation}\label{eq:Model-talenti}
\left\{ \begin{array}{rl}
-\Delta u + (-\Delta)^s u + u^q = \dfrac{f(x)}{u^{\gamma}} & \text{ in } \Omega,\\
u>0 & \text{ in } \Omega,\\
u=0 & \text{ on } \mathbb{R}^{n}\setminus \Omega.
\end{array}\right.
\end{equation}
Then
\begin{equation}\label{th:talenti2}
u^\ast(\tau) \leq\left((\gamma+1)v^\ast(\tau)\right)^{\frac{1}{\gamma+1}}, \quad \text{{\color{red}} for a.e. }\tau \in(0,|\Omega|),
\end{equation}
where $v(x)=v^\ast\left(\omega_n|x|^n\right)$ is the unique radial
solution to the problem
\begin{equation}\label{eq:comp-talenti}
\begin{cases}-\Delta v=f^{\sharp} & \text { in } \Omega^{\sharp} \\ v=0 & \text { on } \partial\Omega^{\sharp} .\end{cases}
\end{equation}
\end{theorem}

\begin{proof}
We split the proof into different steps. 
  \vspace{0.2cm}

-\,\,\textsc{Step I): Approximating problems.} For every $k\in \mathbb{N}$ and $x\in\Omega$ we consider the following nonsingular approximating problems 
%, let us define 
%\begin{equation}\label{def-fk-gk}
%f_k(x):=T_k(f(x))=\min\left\{f(x),k\right\}\quad \text{and}\quad g_k(x):=T_k(u^q(x))=\min\left\{u^q(x),k\right\}.
%\end{equation}
\begin{equation}\label{pk}
\left\{
\begin{array}{ll}
-\Delta u_k+(-\Delta)^s u_k +g_k(u_k)=\dfrac{f_k}{(u_k+\frac{1}{k})^{\gamma}}\qquad& \mbox{in }\Omega
\\
u_k>0& \mbox{in }\Omega
\\
u_k=0& \mbox{on }\mathbb{R}^n\setminus\Omega,
\end{array}
\right.
\end{equation}
where, recalling \eqref{eq:def_Tk} and exploiting that $f\geq 0$, we have set
\begin{align*}
	\mathrm{i)}\,\,&f_k= T_{k}(f) = \min\{f,k\}; \\
	\mathrm{ii)}\,\,&
	g_{k}(s)= \begin{cases}
		T_{k}(g(s)) & \text{if $s\geq 0$}\\
		0 & \text{if $s < 0$}
	\end{cases} 
	\,\,\,= \begin{cases}
		\min\{s^p,k\} & \text{if $s\geq 0$} \\
		0 & \text{if $s < 0$}
	\end{cases}
\end{align*}
By Lemma \ref{eq:Approx_Problem}, for every $k \in \N$ problem \eqref{pk} has a unique nonnegative solution belonging to $\XX^{1,2}(\Omega)\cap L^{\infty}(\Omega)$. Setting $\tilde u_k:=u_k+\frac{1}{k}$, problem \eqref{pk} reads as
\begin{equation*}
\begin{cases}-\Delta \tilde{u}_{k}+(-\Delta)^s\tilde u_k+g_k\left(\tilde u_k-\frac 1 k\right)=\dfrac{f_k}{\tilde{u}_{k}^{\gamma}} & \text { in } \Omega  \\ \tilde{u}_{k}>\frac{1}{k} & \text { in } \Omega \\ \tilde{u}_{k}=\frac{1}{k} & \text { on } \R^n\setminus \Omega .\end{cases}
\end{equation*}
In particular,
\begin{equation}\label{weakapprox}
\begin{split}
\int_\Omega\nabla \tilde u_k(x)\cdot\nabla\varphi(x)\,dx&+\iint_{\R^{2n}}\frac{(\tilde u_k(x)-\tilde u_k(y))(\varphi(x)-\varphi(y))}{|x-y|^{n+2s}}\,dx\,dy\\&+\int_\Omega g_k\left(\tilde u_k(x)-\frac 1 k\right)\varphi(x) \,dx
=\int_{\Omega}\frac{f_k(x)}{\tilde u_k(x)^{\gamma}}\,\varphi(x)\,dx
\end{split}
\end{equation}
for every $\varphi\in \XX^{1,2}(\Omega)$.
%Moreover, by Lemma \ref{lem:propunModelCase}, the sequence $u_k$ is non-deceasing, $u_k>0$ in $\Omega$, and, for every subset $\omega \subset\subset \Omega$, there exists a positive constant $c_\omega$, independent of $k$, such that $u_k (x)\ge c_\omega>0$ for almost every $x \in \omega$ and $k \in \N$.

\vspace{0.2cm}

%\noindent \emph{Step 2. Reduction to the radial case}.
Let us fix $\frac{1}{k}< t<\|u_k\|_{L^\infty(\Omega)}$ and $h>0$. We consider the following test function 
$$\varphi(x)=\mathcal{G}_{t,h}(\tilde u_k(x)),$$ 
where $\mathcal{G}_{t,h}(\theta)$ is defined by
\begin{equation}\label{def-G}
\mathcal{G}_{t,h}(\theta)=\begin{cases}
h \qquad & \mbox{if } \theta>t+h\\
\theta-t & \mbox{if } t<\theta\leq t+h\\
0  & \mbox{if } \theta\leq t.
\end{cases}
\end{equation}
We note that $\mathcal{G}_{t,h}(\tilde u_k)\in \XX^{1,2}(\Omega)$, so we can use it the weak formulation of solution \eqref{weakapprox}, obtaining
\begin{equation}\label{b}
\begin{aligned}
\int_{\Omega}\nabla \tilde u_k(x)&\cdot\nabla\mathcal{G}_{t,h}(\tilde u_k(x))\,dx\\
&+\iint_{\R^{2n}}\frac{(\tilde u_k(x)-\tilde u_k(y))\left(\mathcal{G}_{t,h}(\tilde u_k(x))-\mathcal{G}_{t,h}(\tilde u_k(y))\right)}{|x-y|^{n+2s}}\,dx\,dy\\
&+\int_\Omega
g_k\left(\tilde u_k(x)-\frac 1 k \right)\mathcal{G}_{t,h}(\tilde u_k(x))\,dx\\
&=\int_{\Omega}\frac{f_k(x)}{\tilde u_k(x)^{\gamma}}\,\mathcal{G}_{t,h}(\tilde u_k(x))\,dx.
\end{aligned}
\end{equation}
We first deal with the left-hand side in \eqref{b}. To this end we observe that, setting\begin{equation*}
\mathcal{H}:= (\tilde u_k(x)-\tilde u_k(y))\left(\mathcal{G}_{t,h}(\tilde u_k(x))-\mathcal{G}_{t,h}(\tilde u_k(y))\right),
\end{equation*}
by the monotonicty of $\mathcal{G}_{t,h}$ we see that
$\mathcal{H}\geq 0$.
Therefore, 
\begin{equation*}
 \iint_{\R^{2n}}\frac{(\tilde u_k(x)-\tilde u_k(y))\left(\mathcal{G}_{t,h}(\tilde u_k(x))-\mathcal{G}_{t,h}(\tilde u_k(y))\right)}{|x-y|^{n+2s}}\,dx\,dy\geq 0.
\end{equation*}
Moreover, since $g_k\left(\tilde u_k-\frac{1}{k}\right)\geq 0$ and $\mathcal{G}_{h,t}(\tilde u_k)\geq 0$, we have also
\begin{equation*}
    \int_\Omega g_k\left(\tilde u_k-\frac{1}{k}\right)\mathcal{G}_{t,h}(\tilde u_k(x))\,dx\geq 0,
\end{equation*}
thus
\begin{equation}\label{bb}
\int_\Omega\nabla \tilde u_k(x)\cdot\nabla\mathcal{G}_{t,h}(\tilde u_k(x))\,dx\leq \int_{\Omega}\frac{f_k(x)}{\tilde u_k(x)^{\gamma}}\,\mathcal{G}_{t,h}(\tilde u_k(x))\,dx.
\end{equation}
Dividing by $h>0$ the right-hand side of \eqref{bb} and using \eqref{def-G}, we obtain
\begin{equation*}
\begin{aligned}
%\dfrac{1}{h} \int_{\Omega}\frac{f_k(x)}{\left(u_k(x)+\frac{1}{k}\right)^{\gamma}}&\,\mathcal{G}_{t,h}(u_k(x))\,dx \leq \|f\|_{L^\infty(\Omega)}\int_{\{u_k>t+h\}}\frac{1}{\left(u_k(x)+\frac{1}{k}\right)^\gamma}\,dx\\
%&+\|f\|_{L^\infty(\Omega)}\int_{\{t<u_k\leq t+h\}}\frac{1}{\left(u_k(x)+\frac{1}{k}\right)^\gamma}\left( \frac{u(x)-t}{h}\right)\,dx
\dfrac{1}{h} \int_{\Omega}\dfrac{f_k(x)}{\tilde u_k(x)^{\gamma}}\,\mathcal{G}_{t,h}(\tilde u_k(x))\,dx = &\int_{\{\tilde u_k>t+h\}}\dfrac{f_k(x)}{\tilde u_k(x)^{\gamma}}\,dx\\
&\ \ \ \ +\int_{\{t<\tilde u_k\leq t+h\}}\dfrac{f_k(x)}{\tilde u_k(x)^{\gamma}}\left( \frac{\tilde u(x)-t}{h}\right)\,dx\end{aligned}
\end{equation*}
which implies, %setting $\tilde{u}_k=u_k+\frac 1k$,
\begin{equation}\label{lim-RHS}
	\begin{aligned}
    \lim_{h\to 0^+}\dfrac{1}{h}& \int_{\Omega}\dfrac{f_k(x)}{\tilde u_k(x)^{\gamma}}\,\mathcal{G}_{t,h}(u_k(x))\,dx =\int_{\{\tilde u_k>t\}}\dfrac{f_k(x)}{\tilde u_k(x)^{\gamma}}\,dx\\
    %&\leq \int_{\{\tilde u_k>t\}}\dfrac{f_k(x)}{\tilde u_k(x)^\gamma}\,dx.
\end{aligned}
\end{equation}
While for the left-hand side of \eqref{bb}, we have
\begin{equation*}
    \int_\Omega\nabla \tilde u_k(x)\cdot\nabla\mathcal{G}_{t,h}(\tilde u_k(x))\,dx=\int_{\{t<\tilde u_k\leq t+h\}}|\nabla \tilde u_k(x)|^2 \,dx
\end{equation*}
and dividing by $h$ and letting $h\to 0$, we get
\begin{equation}\label{lim-LHS}
    \lim_{h\to0}\frac{1}{h}\int_{\{t<\tilde u_k\leq t+h\}}|\nabla \tilde u_k(x)|^2\, dx=-\frac{d}{dt}\int_{\{\tilde u_k>t\}}|\nabla \tilde u_k(x)|^2\, dx.
\end{equation}
Now by applying the Coarea formula we get
$$
\int_{\{\tilde u_k>t\}}|\nabla \tilde u_k(x)| \,dx=\int_t^{\infty} P(\{x \in \Omega: \tilde u_k(x)>r\}) \,d r, \quad \text{for } t>\frac{1}{k}
$$
where $P(\cdot)$ denotes the perimeter. Then, by the isoperimetric inequality, we have
\begin{equation}\label{use-isop}
n \,\omega_n^{1/n}\, \mu_{\tilde u_k}(t)^{1-1/n} \leq P(\{x \in \Omega: \tilde u_k>t\})=-\frac{d}{d t} \int_{\{\tilde u_k>t\}}|\nabla \tilde u_k(x)| \,dx .
\end{equation}
Moreover, H\"older's inequality gives
\begin{equation}\label{use-holder}
\begin{aligned}
-\dfrac{d}{d t}& \int_{\{\tilde u_k>t\}}|\nabla \tilde u(x)| \,dx  =\lim _{h \rightarrow 0} \dfrac{1}{h} \int_{\{t<\tilde u_k \leq t+h\}}|\nabla \tilde u_k(x)| \,dx \\
& \leq \lim _{h \rightarrow 0}\left(\dfrac{1}{h} \int_{\{t<\tilde u_k \leq t+h\}}|\nabla \tilde u_k(x)|^2 \,dx\right)^{1 /2}\left(\dfrac{\mu_{\tilde u_k}(t)-\mu_{\tilde u_k}(t+h)}{h}\right)^{1/2} \\
& =\left(-\dfrac{d}{d t} \int_{\{\tilde u_k>t\}}|\nabla \tilde u_k(x)|^2 \,d x\right)^{1 / 2} \left(-\mu_{\tilde u_k}^{\prime}(t)\right)^{1 /2}
\end{aligned}
\end{equation}
Thus, combing \eqref{use-isop} and \eqref{use-holder}, we obtain
\begin{equation}\label{deriv-mu}
n^2\omega_n^{2/n} \dfrac{\mu_{\tilde u_k}(t)^{2-2/n}}{-\mu_{\tilde u_k}^{\prime}(t)} \leq -\frac{d}{d t} \int_{\{\tilde u_k>t\}}|\nabla \tilde u_k(x)|^2 \,d x.
\end{equation}
Finally, from \eqref{bb}, \eqref{lim-RHS}, \eqref{lim-LHS} and \eqref{deriv-mu}, we deduce
\begin{equation}\label{ineq-u-k}
    n^2\omega_n^{2/n} \dfrac{\mu_{\tilde u_k}(t)^{2-2/n}}{-\mu_{\tilde u_k}^{\prime}(t)} \leq \int_{\{\tilde  u_k>t\}}\dfrac{f_k(x)}{\tilde u_k(x)^\gamma}\,dx\quad \text{for a.e. }t\in\Big(\frac{1}{k},\|u\|_{L^\infty(\Omega)}\Big)
\end{equation}
Then, recalling Hardy-Littlewood inequality and 
that facts that $f_k \leq f$ and $t>\frac 1k$, we have
$$
n^2 \omega_n^{2 / n} \frac{\mu_{\tilde{u}_k}(t)^{2-2 / n}}{\left(-\mu_{\tilde{u}_k}\right)^{\prime}(t)} %\leq \int_{\tilde{u}_k(x)>t} \frac{f_k(x)}{\tilde{u}_k(x)^\gamma} \,d x
<\frac{1}{t^\gamma} \int_0^{\mu_{\tilde{u}_k}(t)} f_k^\ast(\sigma) \,d \sigma 
\leq \frac{1}{t^\gamma} \int_0^{\mu_{\tilde{u}_k}(t)} f^\ast(\sigma) \,d \sigma
$$
which can be rewritten as 
\begin{equation}\label{bb2-new}
\tilde{u}_k^\ast(\tau)^\gamma\left(-\tilde{u}_k^\ast(\tau)^{\prime}\right) \leq \frac{1}{n^2 \omega_n^{2 / n}} \,\tau^{-2+2 / n} \int_0^\tau f^\ast(\sigma) \,d\sigma, \quad \text { for a.e. } \tau \in(0,|\Omega|).
\end{equation}
Integrating \eqref{bb2-new} on the interval $(\tau,|\Omega|)$, we get
\begin{equation}\label{stima-tilde-u-k}
\tilde{u}_k^\ast(\tau)^{\gamma+1} \leq \frac{\gamma+1}{n^2 \omega_n^{2 / n}} \int_\tau^{|\Omega|} r^{-2+2 / n}\left(\int_0^r f^\ast(\sigma) \,d \sigma\right) \,d r, \quad\text{for a.e. } \tau\in(0,|\Omega|).
\end{equation}

-\,\,\textsc{Step 2): Local symmetrized nonsingular problems.} Let $v$ be the solution to problem \eqref{eq:comp-talenti}. Due to the radial symmetry of $v$, we have 
$$v(x)=v^\sharp(x)=v^*(\omega_n|x|^n).$$
So, arguing as in \textsc{Step 1)}, we find
\begin{equation*}
n^2 \omega_n^{2 / n} \dfrac{\mu_{v}(t)^{2-2/n}}{-\mu_{v}^{\prime}(t)} =\int_{\{v>t\}} f^\sharp(x)\,d x=\int_0^{\mu_v(t)}f^\ast(\sigma)\, d\sigma
\end{equation*}
which can be rewritten as
$$
-v^\ast(\tau)^{\prime} =\frac{1}{n^2 \omega_n^{2 / n}} \tau^{-2+2 / n} \int_0^\tau f^\ast(\sigma) \,d\sigma, \quad \text {for a.e. } \tau \in(0,|\Omega|)
$$
Again, integrating \eqref{bb2-new} on the interval $(\tau,|\Omega|)$, we get
\begin{equation}\label{eq-v}
v^\ast(\tau)= \frac{1}{n^2 \omega_n^{2 / n}} \int_\tau^{|\Omega|} r^{-2+2 / n}\left(\int_0^r f^\ast(\sigma) \,d \sigma\right) \,d r, \quad\text{for a.e. } \tau\in(0,|\Omega|).
\end{equation}
which together \eqref{stima-tilde-u-k} ensures that 
\begin{equation}\label{stima-punt-u-k}
    \tilde u_k^\ast(\tau)^{\gamma+1}\leq (\gamma+1)\,v^\ast(\tau)\quad \text{ for a.e. }\tau \in(0,|\Omega|).
\end{equation}

-\,\,\textsc{Step 3): Passing to the limit as $k \to +\infty$.} By Lemma \ref{lem:First_Approx}, we know that the sequences $u_k$ is bounded in $\XX^{1,2}(\Omega)$. Hence, up to subsequences, $\tilde u_k$ converges to $\tilde u \in \XX^{1,2}(\Omega)$, weakly in $\XX^{1,2}(\Omega)$, strongly in $L^r(\Omega)$ for any $r\in [1,2^\ast)$ and a.e. in $\Omega$, where $u$, is a weak solution to problems \eqref{eq:Model-talenti}. Thus by property iv) of rearrangement, we have also that, up to subsequence, $\tilde u_k^\ast$ converges strongly to $u^\ast$ in $L^r(0,|\Omega|)$, for any $r\in[1,2^\ast)$, and a.e. in $(0,|\Omega|).$ So, passing to the limit as $k\to\infty$ in \eqref{stima-punt-u-k}, we obtain \eqref{th:talenti2}. This closes the proof.
\end{proof}

\end{document}
